% Check criterion, deliverable, and the Decision Log question of Day 12.
% The criterion comes from the Day 12 section of the syllabus and from
% Labs/day12/README.md.
\begin{frame}{Check criterion}
  \small
  The instructor checks this at the end of the lab:
  \begin{itemize}
    \item \cmark\ \code{decide.py} prints \code{Task C1: complete}, and
          \code{cloud\_check.py} prints \code{Task D1: complete}.
    \item \cmark\ Link connected: a spoken ``yes'' turns the LED on, ``no''
          turns it off, and the decision appears on the cloud broker.
    \item \cmark\ Link cut: the LED follows the keywords, and the queue of
          the forwarder grows.
    \item \cmark\ After the link returns, \code{cloud\_check.py} prints
          \code{RESULT: PASS - no lost message}.
    \item \cmark\ The Decision Log gives numbers and names one trade-off.
  \end{itemize}
  \vskip 0.4em
  \textbf{Hand in:} the file \code{report.md} with the topic tree and the
  table of \code{cloud\_check.py}.
\end{frame}

\begin{frame}{Decision Log}
  \begin{exerciseblock}
    A cold store has five temperature sensors, a door contact, and a fan.
    The cloud is behind a mobile modem, which loses the link for up to 2
    hours each day.
    \begin{itemize}
      \item Which decisions stay on the device?
      \item What does each device send, how often, with which QoS, and
            which messages are retained?
      \item How large must the queue of the forwarder be, and what happens
            when it is full?
    \end{itemize}
  \end{exerciseblock}
  \vskip 0.2em
  A good Decision Log (about 100 words):
  \begin{itemize}
    \item gives your measured message rate, the queue length during your
          cut, and the time until the queue was empty,
    \item calculates the queue for 2 hours from these numbers,
    \item names what you lose: for example detail of the data against the
          size of the queue.
  \end{itemize}
\end{frame}
